% GENERATED FILE. Edit canonical lessons, then run npm run generate:books.
% canonical-source-hash: fnv1a64:6963d8128268cde7
% canonical-lessons: GE-C349-stellenanzeige, GE-C349-vorstellungsgespraech, GE-C349-faehigkeit, GE-C349-qualifiziert, GE-C349-beherrschen

\chapter{Eine Stelle suchen}
\label{ch:ge-eine-stelle-suchen}
\hlchaptermodality{349}

\begin{chapteropening}
\textbf{By the end of this chapter:} \emph{I can read a job advertisement, talk about a job interview and say what I am able to do.}

Complete the last lesson of chapter 349: I can read a job advertisement, talk about a job interview and say what I am able to do.
\end{chapteropening}

\section[die Stellenanzeige]{die Stellenanzeige — the job advertisement}

\label{lesson:GE-C349-stellenanzeige}



\begin{quote}
Before the new one: say the German for the excuse, then the German for that is why.
\end{quote}

\subsection*{You'll want to know: die Stellenanzeige}
\textbf{die Stellenanzeige} — \textquotedblleft{}the job advertisement\textquotedblright{}.

\begin{quote}
In der Zeitung habe ich eine interessante Stellenanzeige gefunden.
\end{quote}

\begin{quote}
\emph{I found an interesting job advertisement in the newspaper.}
\end{quote}

Plural: \emph{die Stellenanzeigen}. Online it is often \emph{das Stellenangebot}.

\begin{cousinweb}
\emph{Stellenanzeige} joins \emph{Stelle} \textquotedblleft{}place, job\textquotedblright{}, from the family of \emph{stellen} \textquotedblleft{}to put\textquotedblright{} (Proto-Germanic \emph{stallijaną}, EtymDB), and \emph{Anzeige} \textquotedblleft{}notice, advertisement\textquotedblright{}, from \emph{anzeigen} \textquotedblleft{}to show\textquotedblright{}.
\end{cousinweb}

\begin{grammarlens}[title={What you've built}]
One more word for telling, explaining and describing at B1.
\end{grammarlens}

\subsection*{Your turn}
\begin{itemize}
\raggedright
  \item \emph{Say it:} \emph{die Stellenanzeige}
  \item \emph{Say it:} \emph{die Stellenanzeige}, once more
  \item \emph{Say it:} a sentence of your own with \emph{die Stellenanzeige}
  \item \emph{Recall:} say the German for the excuse, then the German for that is why, then say \emph{die Stellenanzeige} again
\end{itemize}

\begin{tcolorbox}[breakable,colback=teal!4,colframe=teal!35!black,title={Before you move on}]
What does die Stellenanzeige mean? (\textquotedblleft{}The job advertisement\textquotedblright{}.) Say it once more.
\end{tcolorbox}
\section[das Vorstellungsgespräch]{das Vorstellungsgespräch — the job interview}

\label{lesson:GE-C349-vorstellungsgespraech}



\begin{quote}
Before the new one: say the German for that is why, then the German for the job advertisement.
\end{quote}

\subsection*{You'll want to know: das Vorstellungsgespräch}
\textbf{das Vorstellungsgespräch} — \textquotedblleft{}the job interview\textquotedblright{}.

\begin{quote}
Morgen habe ich ein Vorstellungsgespräch bei einer Bank.
\end{quote}

\begin{quote}
\emph{Tomorrow I have a job interview at a bank.}
\end{quote}

Plural: \emph{die Vorstellungsgespräche}. \emph{Sich vorstellen} — to introduce yourself — is at its heart.

\begin{cousinweb}
\emph{Vorstellungsgespräch} joins \emph{Vorstellung} \textquotedblleft{}introduction\textquotedblright{}, from \emph{vorstellen}, built on \emph{stellen} \textquotedblleft{}to put\textquotedblright{} (EtymDB), and \emph{Gespräch} \textquotedblleft{}conversation\textquotedblright{}, from Old High German \emph{gisprāhhi} (EtymDB): the talk where you introduce yourself.
\end{cousinweb}

\begin{grammarlens}[title={What you've built}]
One more word for telling, explaining and describing at B1.
\end{grammarlens}

\subsection*{Your turn}
\begin{itemize}
\raggedright
  \item \emph{Say it:} \emph{das Vorstellungsgespräch}
  \item \emph{Say it:} \emph{das Vorstellungsgespräch}, once more
  \item \emph{Say it:} a sentence of your own with \emph{das Vorstellungsgespräch}
  \item \emph{Recall:} say the German for that is why, then the German for the job advertisement, then say \emph{das Vorstellungsgespräch} again
\end{itemize}

\begin{tcolorbox}[breakable,colback=teal!4,colframe=teal!35!black,title={Before you move on}]
What does das Vorstellungsgespräch mean? (\textquotedblleft{}The job interview\textquotedblright{}.) Say it once more.
\end{tcolorbox}
\section[die Fähigkeit]{die Fähigkeit — the ability, the skill}

\label{lesson:GE-C349-faehigkeit}



\begin{quote}
Before the new one: say the German for the job advertisement, then the German for the job interview.
\end{quote}

\subsection*{You'll want to know: die Fähigkeit}
\textbf{die Fähigkeit} — \textquotedblleft{}the ability, the skill\textquotedblright{}.

\begin{quote}
Für diese Stelle braucht man die Fähigkeit, im Team zu arbeiten.
\end{quote}

\begin{quote}
\emph{For this job you need the ability to work in a team.}
\end{quote}

Plural: \emph{die Fähigkeiten}. From \emph{fähig} — able, capable.

\begin{cousinweb}
\emph{Fähigkeit} is built on \emph{fähig} \textquotedblleft{}able\textquotedblright{}, inherited from Middle High German \emph{gevæhic}, from \emph{vāhen} \textquotedblleft{}to catch, to seize\textquotedblright{} (EtymDB): someone who can take hold of a task.
\end{cousinweb}

\begin{grammarlens}[title={What you've built}]
One more word for telling, explaining and describing at B1.
\end{grammarlens}

\subsection*{Your turn}
\begin{itemize}
\raggedright
  \item \emph{Say it:} \emph{die Fähigkeit}
  \item \emph{Say it:} \emph{die Fähigkeit}, once more
  \item \emph{Say it:} a sentence of your own with \emph{die Fähigkeit}
  \item \emph{Recall:} say the German for the job advertisement, then the German for the job interview, then say \emph{die Fähigkeit} again
\end{itemize}

\begin{tcolorbox}[breakable,colback=teal!4,colframe=teal!35!black,title={Before you move on}]
What does die Fähigkeit mean? (\textquotedblleft{}The ability, the skill\textquotedblright{}.) Say it once more.
\end{tcolorbox}
\section[qualifiziert]{qualifiziert — qualified}

\label{lesson:GE-C349-qualifiziert}



\begin{quote}
Before the new one: say the German for the job interview, then the German for the ability.
\end{quote}

\subsection*{You'll want to know: qualifiziert}
\textbf{qualifiziert} — \textquotedblleft{}qualified\textquotedblright{}.

\begin{quote}
Sie ist für die Stelle sehr gut qualifiziert.
\end{quote}

\begin{quote}
\emph{She is very well qualified for the job.}
\end{quote}

The participle of \emph{qualifizieren}; \emph{sich qualifizieren} — to qualify, also in sport.

\begin{cousinweb}
\emph{qualifiziert} is the participle of \emph{qualifizieren}, a learned borrowing from Medieval Latin \emph{quālificāre} \textquotedblleft{}to give a quality to\textquotedblright{}, built on Latin \emph{quālis} \textquotedblleft{}of what kind\textquotedblright{} (Lewis \& Short).
\end{cousinweb}

\begin{grammarlens}[title={What you've built}]
One more word for telling, explaining and describing at B1.
\end{grammarlens}

\subsection*{Your turn}
\begin{itemize}
\raggedright
  \item \emph{Say it:} \emph{qualifiziert}
  \item \emph{Say it:} \emph{qualifiziert}, once more
  \item \emph{Say it:} a sentence of your own with \emph{qualifiziert}
  \item \emph{Recall:} say the German for the job interview, then the German for the ability, then say \emph{qualifiziert} again
\end{itemize}

\begin{tcolorbox}[breakable,colback=teal!4,colframe=teal!35!black,title={Before you move on}]
What does qualifiziert mean? (\textquotedblleft{}Qualified\textquotedblright{}.) Say it once more.
\end{tcolorbox}
\section[beherrschen]{beherrschen — to master, to have a good command of; to control}

\label{lesson:GE-C349-beherrschen}



\begin{quote}
Before the new one: say the German for the ability, then the German for qualified.
\end{quote}

\subsection*{You'll want to know: beherrschen}
\textbf{beherrschen} — \textquotedblleft{}to master, to have a good command of; to control\textquotedblright{}.

\begin{quote}
Er beherrscht drei Sprachen.
\end{quote}

\begin{quote}
\emph{He has a good command of three languages.}
\end{quote}

Inseparable: \emph{er beherrscht, er hat beherrscht}. \emph{Sich beherrschen} — to control yourself.

\begin{cousinweb}
\emph{beherrschen} is built with \emph{be-} on \emph{herrschen} \textquotedblleft{}to rule\textquotedblright{}, inherited from Middle High German \emph{hērschen} and Old High German \emph{hērisōn} (EtymDB): you rule over a skill.
\end{cousinweb}

\begin{grammarlens}[title={What you've built}]
One more word for telling, explaining and describing at B1.
\end{grammarlens}

\subsection*{Your turn}
\begin{itemize}
\raggedright
  \item \emph{Say it:} \emph{beherrschen}
  \item \emph{Say it:} \emph{beherrschen}, once more
  \item \emph{Say it:} a sentence of your own with \emph{beherrschen}
  \item \emph{Recall:} say the German for the ability, then the German for qualified, then say \emph{beherrschen} again
\end{itemize}

\begin{tcolorbox}[breakable,colback=teal!4,colframe=teal!35!black,title={Before you move on}]
What does beherrschen mean? (\textquotedblleft{}To master, to have a good command of; to control\textquotedblright{}.) Say it once more.
\end{tcolorbox}
